\documentclass[10pt]{article}

\usepackage[utf8]{inputenc}

\usepackage[T1]{fontenc}

\usepackage{amsmath}

\usepackage{amsfonts}

\usepackage{amssymb}

\usepackage[version=4]{mhchem}

\usepackage{stmaryrd}

\usepackage{graphicx}

\usepackage[export]{adjustbox}

\graphicspath{ {./images/} }



\begin{document}

$\left.Z_{n}\right), 0=(0, \ldots, 0)$. Групповой закон на множествах $H(B)=\operatorname{nil}(B)^{n}$ задается формулами



\[

\left(x_{1}, \ldots, x_{n}\right) \bigcirc\left(y_{1}, \ldots, y_{n}\right)=\left(z_{1}, \ldots, z_{n}\right),

\]



где $z_{i}=F_{i}\left(x_{1}, \ldots, x_{n}, y_{1}, \ldots, y_{n}\right)$; в силу нильпотентности $x$ и $y$ все члены рядов кроме конечного числа равны 0 . И любой формальный групповой закон задает на nil $(B)^{n}$ структуры групп с помощью формул (*) и превращает функтор $D^{n}$ в Ф. Г. Ли. Понятие формального группового закона, и тем самым понятие Ф. г. Ли, обобщается на случай произвольных коммутативных базисных колец (см. [2], [5]). Иногда под Ф. г. понимают лишь Ф. г. Ли или даже формальный групповой закон.



Так же как и для Ли локальных әрупn, определяется алгебра Ли Ф. г. Ли. Над полями $k$ характеристики 0 сопоставление Ф. г. Ли ее алгебры Ли определяет эквивалентность соответствующих категорий. В характеристике $p>0$ ситуация сложнее. Так, над алгебраически замкнутым полем (при $p>0$ ) существует счетное число погарно неизоморфных одномерных коммутативных Ф. г. Ли [1], в то время как все одномерные алгебры Ли изоморфны [3]. Над совершенными полями конечной характеристики коммутативные Ф. г. Ли классифицируются с помощью модулей Дьёдонне (см. $[1,6])$



Теория Ф. г. над полями обобщается на случай произвольных базисных формальных схем [7].



Лит.: [1] М а н и н Ю. И., "Успехи матем. наук", 1963



\includegraphics[max width=\textwidth, center]{2023_07_09_0617a9df68e53b0f4496g-1}

1973; [3] С'е р р Ж.- П., Алгебры Ји и групшы Ји, пер. с англ', и франц., М., 1969; [4] Х а р т с X о р н Р., Алгебраическая геометрия, пер. с англ., M., 1981; [5] L a z a r d M., Commutative formal groups, B.- Hdlb.- N. Y., 1975; [6] F o n t a i n e J. - M., Groupes $p$-divisibles sur les corps locaux, P., 1977, [7] M a zu u r B., T a t e J., B KH.: Arithmetic and geometry, v. 1 , Boston - Básel - Stuttg., 1983, p. 195-237. F. Г. Зархин.



ФОРМАЈЬНАЯ ПРООИЗВОДНАЯ - производная многочлена, рациональной функции или формального степенного ряда, определяемая чисто алгебраически (без использования повятия предельного перехода) и имеющая смысл для любого кольца коэфффициентов. Для многочлена



\[

F(X)=\sum_{i=0}^{n} a_{i} X^{i}

\]



(или степенного ряда



\[

\left.A(X)=\sum_{i=0}^{\infty} b_{i} X^{i}\right)

\]



Ф. п. $F^{\prime}(X)$ определяется как $\sum_{i=1}^{n} i a_{i} X^{i-1}$ (соответственно как $\left.\sum_{i=1}^{\infty} i b_{i} X^{i-1}\right)$, а для рациональной функции $f(X)=P(X) / Q(X)$ - это рациональная функция $f^{\prime}(X)=\left\{P^{\prime}(X) Q(X)-Q^{\prime}(X) P(X)\right\} /[Q(X)]^{2}$.



Аналюгично определяются Ф. П. выспих порядков и qастные Ф. П. для функций от нескольких переменных.



Для Ф. п. остается справедливым ряд свойств обычной производной. Так, если $F^{\prime}(X)=0$, то $F(X)$ - константа из поля коәфффициентов (в случае характеристики 0 ) п равна $G\left(X^{p}\right)$ (в случае характеристики $\left.p\right)$. Если $x_{0}$ - корень многочлена кратности $k$, то $x_{0}$ является корнем производной $F^{\prime}(X)$ кратности $k-1$.



ФОРМАЛЬНАЯ СИСТЕМА, Л. В. Кузьмин. Т И В н а я ванное исчисление, задаваемое правилами образования выражений этого исчисления и правилами построения выводов в этом исчислении. Выражения Ф. с. рассматриваются как чисто формальные комбинации символов; правила вывода определяют, в каких случаях одно формальное выражение $A$ выводится из других формальных выражений $B_{1}, \ldots, B_{n}$. Если $n=0$, то $A$ наз. а к с и о м о й. Выводы представляют собой либо последовательности, либо древовидные фигуры, составленные из формальных выражений согласно правилам вывода. Если в вершинах дерева вывода находятся только аксиомы, то формальное выражение, завершающее вывод, наз. в ы в о д и м ы м в Ф. с.



Наиболее интересны Ф. с., удовлетворяющие требованиям ә ф е к т и в н о с т и языка и понятия вывода. Это означает, что должна иметься эффективная процедура для распознавания того, является ли произвольная последовательность символов выражением Ф. с. или нет. Такому же требованию должно удовлетворять понятие вывода. Понятие выводимого выражения в әффективных Ф. с., вообще говоря, не является әфффективным.



Понятие Ф. с. - одно из центральных в математич. логике, оно обслуживает нужды как самой математич. логики, так и смежных с ней областей математики.



Наиболее важный класс Ф. с.- формальные теории 1-го порядка (см. [4]), формализующие какую-либо область содержательной математики. Исторически әтот класс Ф. с. возник в связи с программой Д. Гильберта (D. Hilbert) обоснования математики (см. Формализм).



Понятия и методы, выработанные математич. логикой при изучении Ф. с., нашли применения в различных областях математики, напр. в теории групп и теории категорий. Ценность понятия Ф. с. определяется плодотворностью способа исследования, основанного на идея $x$, связанных с этим понятием.



См. также Формальный математический анализ.



Лит.: [1] Г и Ј ь б е р т Д., Основания геометрии, пер. с нем., М.- Л., 1948; [2] К Ј и н и С. К., Введение в метаматематику, пер. с англ., М., 1957; [3] Ч е р ч А., Введение в мате-



\begin{center}

\includegraphics[max width=\textwidth]{2023_07_09_0617a9df68e53b0f4496g-1(1)}

\end{center}



ФОРМАЛЬНОЕ ПРОИЗВЕДЕНИЕ Т р И $Г$ О Н $о$ М $е$ Три ческих р я д о в



\[

\sum_{n=-\infty}^{\infty} c_{n} e^{i n x}, \sum_{n=-\infty}^{\infty} \gamma_{n} e^{i n x}

\]



\begin{itemize}

  \item ряд

\end{itemize}



\[

\sum_{n=-\infty}^{\infty} K_{n} e^{i n x}

\]



гцӨ



\[

K_{n}=\sum_{m=-\infty}^{\infty} c_{m} \gamma_{n-m}

\]



Если $c_{n} \underset{|n| \rightarrow \infty}{\longrightarrow} 0, \sum_{n=-\infty}^{\infty}\left|n \gamma_{n}\right|<\infty$ п



сходится к сумме $\lambda(x)$,



\[

\sum_{n=-\infty}^{\infty} \gamma_{n} e^{i n x}

\]



то ряд



\[

\sum_{n=-\infty}^{\infty}\left(K_{n}-\lambda(x) c_{n}\right) e^{i n x}

\]



сходится к нулю равномерно на $[-\pi, \pi]$. Условие



\[

\sum_{n=-\infty}^{\infty}\left|n \gamma_{n}\right|<\infty

\]



выполняется, если, напр.,



\[

\sum_{n=-\infty}^{\infty} \gamma_{n} e^{i n x}

\]



есть ряд Фурье трижды диффференцируемой функции $\lambda(x)$





\end{document}